% Таблицы 1–3 русской версии для официального шаблона «Диалога» (LaTeX). Требуют \usepackage{booktabs}.
% Источник чисел — reports/paper/ru/paper_ru.md (и отчёты, указанные там). Десятичная запятая: {,} в математике не нужна —
% числа набраны текстом. Строки [ЖДЁТ РАЗМЕТКИ] заполняются после человеческой проверки.

\begin{table}[t]
\centering\small
\caption{Корпуса после отбора (дуги $N = 20$). Для сказок кластер бутстрэпа — сюжет (варианты одной сказки) или сборник;
проза в выборке LLM-дуг — той же длины, что сказки (643--3\,179 слов), не больше одного текста на автора.}
\label{tab:corpora}
\begin{tabular}{llrrrr}
\toprule
форма & источник & текстов & авторов & медиана слов & LLM-дуги \\
\midrule
волшебная сказка & Афанасьев & 179 & 120 сюжетов & 1\,292 & 179 \\
сказка (перевод) & fairytalez & 1\,074 & 109 сборников & 1\,473 & 300 \\
рассказ & Проза.ру & 145\,731 & 31\,568 & 1\,311 & 300 \\
новелла & Проза.ру & 5\,863 & 2\,146 & 1\,205 & 300 \\
фанфик-ориджинал & Ficbook & 13\,477 & 8\,221 & 2\,934 & 300 \\
повесть & Проза.ру & 1\,962 & 1\,263 & 18\,248 & 100 (отрывки) \\
роман & Проза.ру & 1\,079 & 710 & 50\,714 & 100 (отрывки) \\
\bottomrule
\end{tabular}
\end{table}

\begin{table}[t]
\centering\small
\caption{Согласие сигналов тональности с консенсусом шести LLM ($\rho$ Спирмена, золото; 95~\% ДИ — бутстрэп по авторам).}
\label{tab:validation}
\begin{tabular}{llll}
\toprule
сигнал & все (400) & проза (340) & сказки Афанасьева (60) \\
\midrule
классификатор & 0,52 [0,44; 0,60] & 0,58 [0,50; 0,67] & 0,35 [0,07; 0,60] \\
RuSentiLex & 0,41 [0,31; 0,50] & 0,44 [0,35; 0,54] & 0,29 [$-$0,01; 0,55] \\
LLM (против консенсуса остальных) & 0,92 & --- & --- \\
люди против консенсуса LLM & \multicolumn{3}{l}{[ЖДЁТ РАЗМЕТКИ]} \\
\bottomrule
\end{tabular}
\end{table}

\begin{table}[t]
\centering\small
\caption{Волшебная сказка против рассказа той же длины (LLM-дуги, $N = 20$; 95~\% ДИ).}
\label{tab:frame}
\begin{tabular}{ll}
\toprule
мера & значение \\
\midrule
начало: $z(1\text{--}2)$ $-$ середина, рассказ $-$ сказка & $+$0,29 [$+$0,13; $+$0,45] \\
то же без первого сегмента: $z(2\text{--}3)$ & $+$0,16 [$+$0,01; $+$0,31] \\
середина: $D$ по позициям 3--19, $p$ & 0,38 \\
конец: $z(18\text{--}19)$ $-$ середина, сказка $-$ рассказ & $+$0,02 [$-$0,15; $+$0,19] \\
конец: $z(20)$ $-$ середина, сказка $-$ рассказ & $+$0,69 [$+$0,43; $+$0,94] \\
$D$ до сказки (все 20 позиций) & 0,77 [0,40; 0,98] \\
доля $D^2$ на позиции 20 / на позициях 1--2 (без 20) & 0,69 / 0,63 \\
сказок с финальной формулой & 51~\% / 66~\% \\
сегмент 20 без формулы: сказка $-$ рассказ & $+$0,54 [$+$0,28; $+$0,79] \\
люди: (конец $-$ середина), сказка $-$ рассказ & [ЖДЁТ РАЗМЕТКИ] \\
\bottomrule
\end{tabular}
\end{table}
